\section{Perplexity}

在语言模型里，perplexity 是最传统的评估方式之一。把语言模型看成序列分布 $p(x)$ 后，perplexity 可以理解为它给某个数据集 $D$ 赋予的平均困惑程度，常写成

\[
\mathrm{PPL}(D) = \left(\frac{1}{p(D)}\right)^{1/|D|}
\]

训练时我们最小化训练集上的 perplexity，评估时自然会想看测试集上的 perplexity。

\subsection{为什么 perplexity 仍然重要}

\begin{itemize}
    \item 它比 downstream accuracy 更平滑，适合做 scaling law。
    \item 它是通用的语言建模目标，也是训练目标本身。
    \item 甚至可以在下游任务上定义条件 perplexity。
\end{itemize}

\begin{knowledgebox}{perplexity 的局限}
如果评估器要信任模型给出的概率分布，那么 perplexity 本身就变成了``要求模型先诚实地报出自己的不确定性''。这比直接看黑盒输出的准确率更脆弱。
\end{knowledgebox}

\subsection{Perplexity 与 bits-per-byte 的转换}

在实际论文中，不同团队报告 perplexity 时使用的 tokenizer 不同，导致数值无法直接比较。为了统一，社区越来越倾向于使用 \textbf{bits-per-byte}（BPB）作为与 tokenizer 无关的度量。

转换关系如下。首先，perplexity 和 bits-per-token 的关系是：
\[
\text{bits-per-token} = \log_2(\mathrm{PPL})
\]
然后，bits-per-byte 需要除以每个 token 平均包含的字节数：
\[
\text{bits-per-byte} = \frac{\log_2(\mathrm{PPL})}{\text{avg bytes per token}}
\]

\begin{table}[H]
\centering
\begin{tabular}{p{3.5cm}p{3.0cm}p{3.0cm}p{3.5cm}}
\toprule
\textbf{Tokenizer} & \textbf{PPL} & \textbf{bits/token} & \textbf{bits/byte (approx)} \\
\midrule
GPT-2 BPE (avg 3.5 B/tok) & 15.0 & 3.91 & 1.12 \\
Llama BPE (avg 4.0 B/tok) & 8.0 & 3.0 & 0.75 \\
字符级 (1 B/tok) & 4.5 & 2.17 & 2.17 \\
\bottomrule
\end{tabular}
\caption{不同 tokenizer 下，同样质量的模型 PPL 差异很大，但 bits-per-byte 更具可比性}
\end{table}

\begin{importantbox}{为什么 bits-per-byte 比 perplexity 更公平}
假设模型 A 用的 tokenizer 把 ``hello'' 编码成 1 个 token，模型 B 编码成 3 个 token。即使两者给出完全相同的字节级概率分布，模型 B 的 perplexity 也会更低（因为它分更多步做预测，每步更容易）。bits-per-byte 通过归一化到字节级别，消除了这种 tokenizer 差异。
\end{importantbox}

一个具体的换算例子：如果某模型的 perplexity 是 10，使用的 tokenizer 平均每个 token 对应 4 个字节，那么：
\[
\text{BPB} = \frac{\log_2(10)}{4} = \frac{3.322}{4} \approx 0.83 \text{ bits/byte}
\]
这个数值的直觉含义是：模型平均需要 0.83 个 bit 来编码原始文本中的每个字节。与信息论中英语文本的经验熵（约 1.0--1.5 bits/byte）相比，这意味着模型已经能相当好地预测文本内容。

\subsection{几种``精神上接近 perplexity''的任务}

课程里提到，某些 cloze 或补全类任务和 perplexity 的思想很接近，比如 LAMBADA、HellaSwag 等。它们本质上仍然是“给定前文，预测后文”，只是形式更贴近下游任务。

\begin{warningbox}{最大化 perplexity 的信念}
一种极端观点是：如果模型分布 $p$ 最终逼近真实分布 $t$，那就能解决所有任务。但这并不意味着对分布的每一部分都同等值得优化，也不意味着它是最有效的通往强模型的路径。
\end{warningbox}

